\begin{center}
    {\LARGE \textbf{2016物理化学}} \\[2ex]
    {\large 时量180分钟 \quad 满分150分}
\end{center}

\vspace{0.5cm}
\noindent\textbf{注意事项:}
\begin{enumerate}[label=\arabic*., parsep=0pt, itemsep=0pt]
    \item 答题前,考生先将自己的姓名、学号、考场号、座位号填写在答题卡上,将条形码准确粘贴在考生信息条形码粘贴区。
    \item 考试结束后,将本试卷与答题纸一并收回。
    \item 回答各个问题时,将答案写在答题纸上,写在本试卷上无效。
\end{enumerate}

\vspace{0.5cm}
\noindent\textbf{一、选择题（30 pts.）}

\begin{enumerate}[label=\arabic*., itemsep=1.5ex]
    \item 某定量均相纯流体从298 K、$10p^\ominus$恒温压缩时，总物系的焓增加，则该物系从298 K、$10p^\ominus$节流膨胀到邻近某一状态时，物系的温度必将
    \begin{tasks}(4)
        \task 升高
        \task 降低
        \task 不变
        \task 不能确定
    \end{tasks}
    \item 两半电池之间使用盐桥，测得电动势为0.059 V，当盐桥拿走，使两溶液接触，这时测得电动势为0.048 V，则液接电势值为
    \begin{tasks}(4)
        \task $0.011\ \mathrm{V}$
        \task $0.011\ \mathrm{V}$
        \task $0.107\ \mathrm{V}$
        \task $-0.107\ \mathrm{V}$
    \end{tasks}
    \item 单原子分子理想气体的$C_{V,\mathrm{m}}=(3/2)R$，温度由$T_1$变到$T_2$时，等压过程体系的熵变$\Delta S_p$与等容过程熵变$\Delta S_V$之比是
    \begin{tasks}(4)
        \task $1:1$
        \task $2:1$
        \task $3:5$
        \task $5:3$
    \end{tasks}
    \item $\ce{Na2CO3}$可形成三种水盐：$\ce{Na2CO3\cdot H2O}$、$\ce{Na2CO3\cdot 7H2O}$及$\ce{Na2CO3\cdot 10H2O}$，常压下将$\ce{Na2CO3(s)}$投入其水溶液中，待达三相平衡时，一相是$\ce{Na2CO3}$水溶液，一相是$\ce{Na2CO3(s)}$，则另一相是
    \begin{tasks}(2)
        \task 冰
        \task $\ce{Na2CO3\cdot 10H2O(s)}$
        \task $\ce{Na2CO3\cdot 7H2O(s)}$
        \task $\ce{Na2CO3\cdot H2O(s)}$
    \end{tasks}
    \item 在应用电位计测定电动势的实验中，通常必须用到
    \begin{tasks}(2)
        \task 标准电池
        \task 标准氢电极
        \task 甘汞电极
        \task 活度为1的电解质溶液
    \end{tasks}
    \item 用铜电极电解$0.1\mathrm{mol\cdot kg^{-1}}$的$\ce{CuCl2}$水溶液，阳极上的反应为
    \begin{tasks}(2)
        \task $\ce{2Cl- -> Cl2 + 2e-}$
        \task $\ce{Cu -> Cu^{2+} + 2e-}$
        \task $\ce{Cu -> Cu+ + e-}$
        \task $\ce{2OH- -> H2O + \frac{1}{2}O2 + 2e-}$
    \end{tasks}
    \item 反应$\ce{2N2O5 -> 4NO2 + O2}$的速率常数单位是$\mathrm{s^{-1}}$。对该反应的下述判断哪个对？
    \begin{tasks}(4)
        \task 单分子反应
        \task 双分子反应
        \task 复合反应
        \task 不能确定
    \end{tasks}
    \item $1\ \mathrm{mol}$某气体的状态方程为$pV_\mathrm{m}=RT+bp$，$b$为不等于零的常数，则下列结论正确的是
    \begin{tasks}(1)
        \task 其焓$H$只是温度$T$的函数
        \task 其内能$U$只是温度$T$的函数
        \task 其内能和焓都只是温度$T$的函数
        \task 其内能和焓不仅与温度$T$有关，还与气体的体积$V_\mathrm{m}$或压力$p$有关
    \end{tasks}
    \item $\ce{CaCl2}$摩尔电导率与其离子的摩尔电导率的关系是
    \begin{tasks}(1)
        \task $\Lambda_\infty(\ce{CaCl2})=\lambda_\mathrm{m}(\ce{Ca^{2+}})+\lambda_\mathrm{m}(\ce{Cl-})$
        \task $\Lambda_\infty(\ce{CaCl2})=\lambda_\mathrm{m}(\ce{Ca^{2+}})+2\lambda_\mathrm{m}(\ce{Cl-})$
        \task $\Lambda_\infty(\ce{CaCl2})=2[\lambda_\mathrm{m}(\ce{Ca^{2+}})+\lambda_\mathrm{m}(\ce{Cl-})]$
        \task $\Lambda_\infty(\ce{CaCl2})=1/2[\lambda_\mathrm{m}(\ce{Ca^{2+}})+\lambda_\mathrm{m}(\ce{Cl-})]$
    \end{tasks}
    \item 下列化学反应在恒压、绝热和只作体积功的条件下进行，体系的温度由$T_1$升高到$T_2$，则此过程的焓变$\Delta H$
    \begin{tasks}(4)
        \task 小于零
        \task 等于零
        \task 大于零
        \task 不能确定
    \end{tasks}
    \item 对于1-1型电解质溶液的摩尔电导率可以看作是正、负离子的摩尔电导率之和，这一规律只适用于
    \begin{tasks}(2)
        \task 强电解质
        \task 弱电解质
        \task 无限稀释电解质溶液
        \task 摩尔浓度为1的溶液
    \end{tasks}
    \item 1 mol 373 K、$p^\ominus$下的水经下列两个不同过程变成373 K、$p^\ominus$下的水汽，(1)等温等压可逆蒸发，(2)真空蒸发。这两个过程中功和热的关系为
    \begin{tasks}(2)
        \task $W_1>W_2$, $Q_1>Q_2$
        \task $W_1<W_2$, $Q_1<Q_2$
        \task $W_1=W_2$, $Q_1=Q_2$
        \task $W_1>W_2$, $Q_1<Q_2$
    \end{tasks}
    \item 将$\ce{AlCl3}$溶于水中全部水解，此体系的组分数$K$是
    \begin{tasks}(4)
        \task 1
        \task 2
        \task 3
        \task 4
    \end{tasks}
    \item 在等温、等压下，电池以可逆方式对外作电功的热效应$Q_R$等于
    \begin{tasks}(4)
        \task $\Delta H$
        \task $zFT(\partial E/\partial T)_p$
        \task $zFE(\partial E/\partial T)_p$
        \task $nEF$
    \end{tasks}
    \item 有一$\ce{ZnCl2}$水溶液，$b=0.002\ \mathrm{mol\cdot kg^{-1}}$，$\gamma_\pm=0.83$，则$a_\pm$为
    \begin{tasks}(4)
        \task $1.66\times10^{-3}$
        \task $2.35\times10^{-3}$
        \task $2.64\times10^{-3}$
        \task $2.09\times10^{-4}$
    \end{tasks}
\end{enumerate}

\vspace{0.5cm}
\noindent\textbf{二、计算题（20 pts.）}

$4\ \mathrm{g}\ \ce{Ar}$（可视为理想气体，其摩尔质量$M(\ce{Ar})=39.95\ \mathrm{g\cdot mol^{-1}}$）在300 K时，压力为506.6 kPa，今在等温下反抗202.6 kPa的恒定外压进行膨胀。试分别求下列两种过程的$Q$，$W$，$\Delta U$，$\Delta H$，$\Delta A$和$\Delta G$。
\begin{enumerate}[label=(\arabic*), itemsep=1ex]
    \item 若变化为可逆过程；
    \item 若变化为不可逆过程。
\end{enumerate}

\vspace{0.5cm}
\noindent\textbf{三、计算题（15 pts.）}

293.15 K时，1 mol苯（$C_{p,\mathrm{m}}(\ce{B})=83.7\ \mathrm{J\cdot K^{-1}\cdot mol^{-1}}$）绝热地与313.15 K的2 mol甲苯（$C_{p,\mathrm{m}}(\ce{T})=96\ \mathrm{J\cdot K^{-1}\cdot mol^{-1}}$）混合，形成理想溶液。计算总熵变。

\vspace{0.5cm}
\noindent\textbf{四、计算题（15 pts.）}

气态正戊烷和异戊烷的$\Delta_\mathrm{f}G_\mathrm{m}^\ominus(298.15\ \mathrm{K})$分别为$-194.4\ \mathrm{kJ\cdot mol^{-1}}$，$-200.8\ \mathrm{kJ\cdot mol^{-1}}$，液体的饱和蒸气压分别为

正戊烷：$\lg(p_i^\ominus/\mathrm{Pa})=3.9714-1065/[(T/\mathrm{K})-41]$

异戊烷：$\lg(p_2^\ominus/\mathrm{Pa})=3.9089-1024/[(T/\mathrm{K})-41]$

计算$25^\circ\mathrm{C}$时异构化反应：正戊烷$\rightleftharpoons$异戊烷在气相中的$K_p$和在液相中的$K_x$，假定气相为理想气体混合物，液相为理想溶液。

\vspace{0.5cm}
\noindent\textbf{五、计算题（15 pts.）}

在$25^\circ\mathrm{C}$时，反应$\ce{H2(g) + Ag2O(s) = 2Ag(s) + H2O(l)}$的恒容热效应$Q_V=-252.79\ \mathrm{kJ\cdot mol^{-1}}$，在$p^\ominus$、298 K下，将上述反应体系构成一可逆原电池，测其电动势的温度系数为$-5.044\times10^{-4}\ \mathrm{V\cdot K^{-1}}$，求$\ce{Ag|Ag2O|OH-}$电极的$\varphi^\ominus$。已知$25^\circ\mathrm{C}$时，水的离子积$K_\mathrm{w}=1\times10^{-14}$。

\vspace{0.5cm}
\noindent\textbf{六、计算题（15 pts.）}

在671～768 K之间，$\ce{C2H5Cl}$气相分解反应（$\ce{C2H5Cl -> C2H4 + HCl}$）为一级反应，速率常数$k(\mathrm{s^{-1}})$和温度$T$的关系式为
\[
\lg(k/\mathrm{s^{-1}}) = -13290/(T/\mathrm{K}) + 14.6
\]
\begin{enumerate}[label=(\arabic*), itemsep=1ex]
    \item 求$E_\mathrm{a}$和$A$；
    \item 在700 K时，将$\ce{C2H5Cl}$通入一个反应器中（$\ce{C2H5Cl}$的起始压力为26664.5 Pa），反应开始后，反应器中压力增大，问需多少时间反应器中压力变为46662.8 Pa？
\end{enumerate}

\vspace{0.5cm}
\noindent\textbf{七、计算题（20 pts.）}

电池$\ce{Cd(s) | Cd(OH)2(s) | NaOH(0.01\ mol\cdot kg^{-1}) | H2(p^\ominus), Pt}$在298 K时电动势$E=0.000\ \mathrm{V}$，$(\partial E/\partial T)_p=0.002\ \mathrm{V\cdot K^{-1}}$，$\varphi^\ominus(\ce{Cd^{2+}|Cd})=-0.403\ \mathrm{V}$。
\begin{enumerate}[label=(\arabic*), itemsep=1ex]
    \item 写出电极反应和电池反应；
    \item 求电池反应的$\Delta_\mathrm{r}G_\mathrm{m}$，$\Delta_\mathrm{r}H_\mathrm{m}$，$\Delta_\mathrm{r}S_\mathrm{m}$；
    \item 求$\ce{Cd(OH)2}$的溶度积常数$K_{\mathrm{sp}}$。
\end{enumerate}